% Source: Corsin Nick/Class Notes/Week 4.pdf, pp. 1--13
\chapter{Week 4 --- The Differential, the Chain Rule \& Taylor Expansions}

\exercisesheet{4}

\textit{Statements quoted verbatim from \texttt{exercises/Ex4\_Analysis2\_eng.pdf} (assigned 9
March 2026, due 16 March 2026). Attribution: Prof. Joaquim Serra, D-MATH, ETH Z\"urich.}

\begin{ainote}[Corsin's recommendations]
Colour key, as given on the page: blue \textbf{= important}, orange \textbf{= semi-important},
red \textbf{= optional}.

\begin{center}
\begin{tabular}{lll}
\textbf{Problem} & \textbf{Priority} & \textbf{Corsin's note} \\
\hline
4.1 & important & ``Use the chain rule'' \\
4.2 & important & ``Recall that path connected $\implies$ connected'' \\
4.3 & semi-important --- parts 1, 2, 3, 5; & ``For 3., use Cauchy--Schwarz in $\mathbb{R}^n$ \\
    & \textbf{parts 4 and 6 important} & applied to convenient vectors.'' \\
4.4 & optional & --- \\
4.5 & optional & --- \\
4.6 & important --- parts 1, 2, 3 & --- \\
\end{tabular}
\end{center}
\end{ainote}

\begin{ainote}
Only the exercises tagged \textbf{important} above are reproduced below. The full sheet
(4.1--4.6, including 4.4 and 4.5) is transcribed in
\texttt{content/exercise-sheets/week-04.tex}.
\end{ainote}

\begin{exercise}[4.1 --- Derivative computation \textnormal{(important)}]
\label{ex:4.1}
Consider the function $u : (x,y) \mapsto x^{\sin(y)}$, defined for
$(x,y) \in (0,\infty)\times\mathbb{R} \subset \mathbb{R}^2$. Compute $\partial_x u$ and
$\partial_y u$.
\end{exercise}

\begin{exercise}[4.2 --- Connected graphs \textnormal{(important)}]
\label{ex:4.2}
Let $U \subset \mathbb{R}^n$ be open and connected and let $f \in C^1(U, \mathbb{R}^m)$. Show that
its graph
$$\Gamma_f := \{(x, f(x)) : x \in U\}$$
is a connected subset of $\mathbb{R}^n \times \mathbb{R}^m$.
\end{exercise}

\begin{exercise}[4.3 --- $p$-norms \textnormal{(semi-important; parts 4 and 6 important)}]
\label{ex:4.3}
For $p \geq 1$ and $x \in \mathbb{R}^n$ define the \textbf{$p$-norm} of $x$ as
$$|x|_p := \Big(\sum_{i=1}^{n} |x_i|^p\Big)^{1/p}.$$
\begin{enumerate}[label=\textbf{\arabic*.}]
  \item For $n = 2$ and $p = 1, 2, 10$ sketch the sets $\{x \in \mathbb{R}^2 : |x|_p \leq 1\}$.
  \item For a given $x \in \mathbb{R}^n$, compute the limit
    $|x|_\infty := \lim_{p\to\infty}|x|_p$.
  \item Using an appropriate inequality that you have seen in class, prove that
    $$\Big(\sum_{i=1}^{n} a_i^{p-1}b_i\Big)^2 \leq \Big(\sum_{i=1}^{n}a_i^p\Big)\Big(\sum_{i=1}^{n}a_i^{p-2}b_i^2\Big),$$
    whenever $a_i, b_i$ are $n$-tuples of positive numbers.
  \item Fix $x, y \in \mathbb{R}^n$ and consider the function $f : [0,1] \to \mathbb{R}$ defined
    as
    \begin{equation*}
    f(t) := |tx + (1-t)y|_p = \Big(\sum_{i=1}^{n}|tx_i + (1-t)y_i|^p\Big)^{1/p}. \tag{1}
    \end{equation*}
    Show that $f$ is convex. You may assume that the coordinates of $x$ and $y$ are all strictly
    positive and use the inequality of the previous point.
  \item Deduce from the previous point that the triangular inequality holds, i.e.\
    $|x+y|_p \leq |x|_p + |y|_p$ for all $x,y \in \mathbb{R}^n$.
  \item What happens for $p \in (0,1)$?
\end{enumerate}
\textit{Official hints:} \textbf{4.3.3} --- use Cauchy--Schwarz and the fact that
$a_i^{p-1}b_i = a_i^{p/2}\cdot a_i^{(p-2)/2}b_i$. \textbf{4.3.4} --- show that $f''(t) \geq 0$; it
is not the simplest derivative, but if you get it right the inequality in 4.3.3 will be just what
you need. \textbf{4.3.5} --- write the convexity inequality between $x$, $y$ and the middle point
$(x+y)/2$. To get the general case from the one with positive coordinates, observe that for
$x \in \mathbb{R}^n$, denoting $\hat x := (|x_1|,\dots,|x_n|)$, we have
$|x+y|_p \leq |\hat x + \hat y|_p$ and $|\hat x|_p = |x|_p$ for all $x,y \in \mathbb{R}^n$.
\end{exercise}

\begin{exercise}[4.6 --- A directional derivative vanishes \textnormal{(important)}]
\label{ex:4.6}
Let $u \in C^1(\mathbb{R}^n)$ and $\nu \in \mathbb{R}^n$. Show that
\begin{enumerate}[label=\textbf{\arabic*.}]
  \item If $\partial_1 u \equiv 0$ then ``$u$ does not depend on $x_1$''; more rigorously: there
    exists a unique function $v \in C^1(\mathbb{R}^{n-1})$ such that
    \begin{equation*}
    u(x_1,\dots,x_n) = v(x_2,\dots,x_n) \quad \text{for all } x \in \mathbb{R}^n. \tag{2}
    \end{equation*}
  \item If $\partial_\nu u \equiv 0$ and $\nu\cdot e_1 \neq 0$ then ``$u$ is a function of $n-1$
    variables''; more rigorously: there exists a unique function $w \in C^1(\mathbb{R}^{n-1})$
    such that
    $$u(x_1,\dots,x_n) = w\!\left(x_2 - \frac{x_1\nu_2}{\nu_1}, \dots, x_n - \frac{x_1\nu_n}{\nu_1}\right) \quad \text{for all } x \in \mathbb{R}^n.$$
  \item \textbf{(*)} What can we conclude if we assume only that $\partial_1 u = 0$ in an open
    connected subset $U \subset \mathbb{R}^n$?
\end{enumerate}
\textit{Official hints:} \textbf{4.6.2} --- apply the mean value theorem to $u(x + t\nu)$.
\textbf{4.6.3} --- consider the domain $U := \{(x,y) : x^2 < y < x^2+1\}$ and the function
$$u(x,y) := \begin{cases} \max\{0, y-2\}^2 & \text{if } x \geq 0, \\ 0 & \text{if } x \leq 0.\end{cases}$$
\end{exercise}

\session{Monday}

\section{The differential}
\label{sec:the_differential}

For $f : \mathbb{R}\to\mathbb{R}$ differentiable at $x_0 \in \mathbb{R}$:
$$f(x_0+h) = f(x_0) + f'(x_0)h + o(h)$$
$$\Updownarrow$$
$$\lim_{h\to 0}\frac{f(x_0+h)-f(x_0)}{h} = f'(x_0)$$
$$\Updownarrow$$
$$\lim_{h\to 0}\left|\frac{f(x_0+h)-f(x_0)-f'(x_0)h}{h}\right| = 0$$

For $F : \mathbb{R}^n \to \mathbb{R}^m$ differentiable at $x_0 \in \mathbb{R}^n$:
$$F(x_0+h) = F(x_0) + \underbrace{DF_{x_0}(h)}_{\text{linear map}} + o(|h|)$$
$$\Updownarrow$$
$$\lim_{h\to 0}\frac{F(x_0+h)-F(x_0)}{|h|} = DF_{x_0}\!\left(\frac{h}{|h|}\right) = \partial_{\frac{h}{|h|}}F(x_0) \quad \text{(directional derivative in direction } h\text{)}$$
$$\Updownarrow$$
$$\lim_{h\to 0}\frac{|F(x_0+h)-F(x_0)-DF_{x_0}(h)|}{|h|} = 0$$

\begin{definition}[Differentiability]
\label{def:differentiable}
We say that $F : \mathbb{R}^n \to \mathbb{R}^m$ is \newterm{differentiable}
(\germanterm{differenzierbar}) at $x_0$ if a linear map $DF_{x_0}$ as above exists, where
$x_0, h \in \mathbb{R}^n$. (Of course, equivalently for $U \subseteq \mathbb{R}^n$ a subset,
$F : U \to \mathbb{R}^m$.)
\end{definition}

\section{The Jacobi matrix}
\label{sec:jacobi_matrix}

Since $DF_{x_0} : \mathbb{R}^n \to \mathbb{R}^m$ is a linear map, there is a \textbf{1--1
correspondence with a matrix once a basis is fixed}! In the standard basis of $\mathbb{R}^n$,
$(e_1,\dots,e_n)$, this is the \newterm{Jacobi matrix} (\germanterm{Jacobi-Matrix}):
$$DF_x \;\cong\; \Jac F(x) = \begin{pmatrix} \dfrac{\partial F_1}{\partial x_1} & \cdots & \dfrac{\partial F_1}{\partial x_n} \\[1ex] \vdots & \ddots & \vdots \\[1ex] \dfrac{\partial F_m}{\partial x_1} & \cdots & \dfrac{\partial F_m}{\partial x_n}\end{pmatrix}$$
($m$ rows, $n$ columns), where we wrote $F(x) = (F_1, \dots, F_m)\transp$ and
$$\frac{\partial F}{\partial x_i} =: \partial_i F(x)$$
is the \newterm{$i$-th partial derivative} (\germanterm{partielle Ableitung}) at $x$:
$$\partial_i F(x) = \lim_{s\to 0}\frac{F(x+se_i)-F(x)}{s}.$$
This defines a function $\partial_i F : U \to \mathbb{R}^m$. If $F$ is differentiable, then
$$\partial_i F(x) = DF_x(e_i).$$

\begin{importantremark}
The partial derivatives can exist even if $F : U \to \mathbb{R}^m$ is not differentiable!
\end{importantremark}

\subsection{Examples}

\begin{example}[Derivative of a curve]
$$\gamma : \mathbb{R}\to\mathbb{R}^2, \qquad t \mapsto \begin{pmatrix}\cos t \\ \sin t\end{pmatrix}$$
What is $\gamma'(\tfrac{\pi}{4}) := D\gamma_{\pi/4}$?
$$\gamma'\!\left(\tfrac{\pi}{4}\right) = \begin{pmatrix}\partial_t\gamma_1(\pi/4) \\ \partial_t\gamma_2(\pi/4)\end{pmatrix} = \begin{pmatrix}-1/\sqrt{2} \\ 1/\sqrt{2}\end{pmatrix}$$
\end{example}

\begin{figurestub}{FIG-W04-01}
Unit circle on axes; the point $\gamma(\pi/4)$ marked with an orange tangent arrow pointing
up-left, and a dotted radius to the point.
\end{figurestub}

This is also called the \textbf{speed of the curve}. We get that
$$\gamma\!\left(\tfrac{\pi}{4}+h\right) = \frac{1}{\sqrt{2}}\begin{pmatrix}1\\1\end{pmatrix} + \frac{h}{\sqrt{2}}\begin{pmatrix}-1\\1\end{pmatrix} + o(|h|).$$

From this example, we see that the \textbf{column} vectors of the Jacobian are the partial
derivatives:
$$\Jac F(x) = \left(\frac{\partial F}{\partial x_1}\ \Big|\ \cdots\ \Big|\ \frac{\partial F}{\partial x_n}\right)$$
and therefore, in the canonical basis, with $v = (v_1,\dots,v_n)\transp$:
$$\partial_v F(x) = DF_x(v) = \Jac F(x)\,v = \sum_{i=1}^{n} v_i\,\frac{\partial F}{\partial x_i}.$$

\begin{example}[Scalar field gradient]
$$f : \mathbb{R}^3\to\mathbb{R}, \qquad (x,y,z) \mapsto 2x^2 + y^2 + 3z^2 - 2xyz$$
What is $\Jac f(x,y,z)$?
$$\Jac f(x,y,z) = (4x - 2yz,\ 2y - 2xz,\ 6z - 2xy) = \left(\frac{\partial f}{\partial x}, \frac{\partial f}{\partial y}, \frac{\partial f}{\partial z}\right) = (\nabla f)\transp$$
\end{example}

From this, we see that the \textbf{rows} of the Jacobian are the \newterm{gradients}
(\germanterm{Gradient}):
$$\Jac F(x) = \begin{pmatrix} \nabla F_1 \\ \hline \vdots \\ \hline \nabla F_m \end{pmatrix}$$

\begin{ainote}
Corsin's sentence reads ``the columns of the Jacobian are the gradients'', but the displayed
matrix stacks $\nabla F_1, \dots, \nabla F_m$ as rows --- which is the correct statement, and
consistent with $(\nabla f)\transp$ just above. Corrected to ``rows''. See \texttt{OQ-09}.
\end{ainote}

\section{Chain rule}
\label{sec:chain_rule}

\begin{theorem}[Chain rule]
\label{thm:chain_rule}
The \newterm{chain rule} (\germanterm{Kettenregel}) states: for
$f : U \subseteq \mathbb{R}^n \to \mathbb{R}^m$, $g : V \subseteq \mathbb{R}^m \to \mathbb{R}^d$
differentiable, it holds:
$$D(g\circ f)(x) = Dg(f(x))\,Df(x)$$
And equivalently for the Jacobians.
\end{theorem}

\begin{exercise}[Jacobian of a composite map]
\label{ex:chain_rule_computation}
Compute $\Jac(g\circ f)$ for
$$f : \mathbb{R}^3\to\mathbb{R}^2,\ x \mapsto (x_1x_2,\ x_2+x_3), \qquad g : \mathbb{R}^2\to\mathbb{R}^2,\ x \mapsto (e^{x_1},\ x_1x_2).$$
\end{exercise}

\begin{exercisesolution}[\cref{ex:chain_rule_computation}]
We have
$$\Jac f(x) = \begin{pmatrix} x_2 & x_1 & 0 \\ 0 & 1 & 1\end{pmatrix}, \qquad \Jac g(x) = \begin{pmatrix} e^{x_1} & 0 \\ x_2 & x_1 \end{pmatrix}.$$
So:
$$
\begin{aligned}
\Jac(g\circ f) &= \begin{pmatrix} e^{x_1x_2} & 0 \\ x_2+x_3 & x_1x_2\end{pmatrix}\begin{pmatrix} x_2 & x_1 & 0 \\ 0 & 1 & 1\end{pmatrix} \\
&= \begin{pmatrix} x_2e^{x_1x_2} & x_1e^{x_1x_2} & 0 \\ x_2^2 + x_2x_3 & 2x_1x_2 + x_1x_3 & x_1x_2 \end{pmatrix}
\end{aligned}
$$
\end{exercisesolution}

Sometimes the \newterm{differential} (\germanterm{das Differential}) is not easily expressed as a
Jacobian. Then it can be convenient to use the formula for \newterm{directional derivatives}
(\germanterm{Richtungsableitung}):
\begin{equation*}
DF_{x_0}(v) = \partial_v F(x_0) = \left.\frac{d}{ds}\right|_{s=0} F(x_0 + sv) \tag{$*$}
\end{equation*}

\begin{importantremark}
The directional derivative $\partial_v F(x_0)$ can exist (depending on $v$ and $x_0$) even if $F$
is not differentiable at $x_0$, i.e.\ $DF_{x_0}$ does not exist. However, if $DF_{x_0}$ exists,
the formula $(*)$ holds.
\end{importantremark}

\begin{exercise}[Differential of $A \mapsto A\transp A$]
\label{ex:transpose_diff}
We identify the inner product space of real-valued matrices $\mathbb{R}^{n\times n}$ with inner
product
$$\langle A, B\rangle = \sum_{i,j=1}^{n} A_{ij}B_{ij} = \Tr(A\transp B)$$
for $A, B \in \mathbb{R}^{n\times n}$, with the euclidean space
$\mathbb{R}^{n^2}\cong\mathbb{R}^{n\times n}$, with the standard inner product. Then, consider
the function
$$F : \mathbb{R}^{n\times n}\to\mathbb{R}^{n\times n}, \qquad A \mapsto A\transp A.$$
Compute the differential at the identity matrix $\id \in \mathbb{R}^{n\times n}$ applied to
$X \in \mathbb{R}^{n\times n}$:
$$DF_{\id}(X) = \partial_X F(\id).$$
\end{exercise}

\begin{exercisesolution}[\cref{ex:transpose_diff}]
We use $(*)$:
$$
\begin{aligned}
DF_{\id}(X) &= \left.\frac{d}{ds}\right|_{s=0} F(\id+sX) \\
&= \left.\frac{d}{ds}\,(\id+sX)\transp(\id+sX)\right|_{s=0} \\
&= \left.\frac{d}{ds}\left[\id + s(X\transp+X) + s^2X\transp X\right]\right|_{s=0} \\
&= X\transp + X
\end{aligned}
$$
\end{exercisesolution}

\session{Friday}

\section{The chain rule along a path}
\label{sec:chain_rule_along_path}

An important special case of the chain rule is that of a composition of a scalar field with a
path. Suppose $f : \mathbb{R}^n \supseteq U \to \mathbb{R}$ and
$\gamma : [0,1]\to\mathbb{R}^n$ are a $C^1$ \newterm{scalar field} (\germanterm{Skalarfeld}) and
path, respectively. Then
$$
\begin{aligned}
D(f\circ\gamma)(t) &= Df(\gamma(t))\cdot D\gamma(t) \\
&= \langle \nabla f(\gamma(t)),\ \gamma'(t)\rangle \\
&= \sum_{i=1}^{n} \frac{\partial f}{\partial \gamma_i}\frac{\partial \gamma_i}{\partial t}
\end{aligned}
$$
where $\gamma = (\gamma_1,\dots,\gamma_n)$ and
$\dfrac{\partial f}{\partial \gamma_i} := \partial_i f(\gamma(t))$.

This operation is so common in physics that it has its own notation: for $f$ as above, we look at
each coordinate as a function of \textbf{time} and differentiate:
$$\frac{d}{dt}f(x_1,\dots,x_n) := Df(x_1(t),\dots,x_n(t)) = \sum_{i=1}^{n}\frac{\partial f}{\partial x_i}\frac{\partial x_i}{\partial t}.$$

\section{Taylor expansions}
\label{sec:taylor_expansions}

In the lecture, you have seen the following formula. For
$f \in C^k(U \subseteq \mathbb{R}^n, \mathbb{R})$, $\bar x \in U$:
$$f(\bar x + h) = \sum_{|\alpha| = 0}^{k} \partial^\alpha f(\bar x)\,\frac{h^\alpha}{\alpha!} + o(|h|^k)$$
where $\alpha$ is a \newterm{multi-index} (\germanterm{Multiindex}) in $\mathbb{N}^n$.

In practice, we never use this formula, but to shed some light on it:

\begin{exercise}[Maximal Taylor expansion in $\mathbb{R}^2$]
\label{ex:taylor_maximal_expansion}
Let $f \in C^3(\mathbb{R}^2,\mathbb{R})$. Write out the maximal Taylor approximation.
\end{exercise}

\begin{exercisesolution}[\cref{ex:taylor_maximal_expansion}]
The multi-indices are:
\begin{itemize}
  \item $|\alpha| = 0$: $\alpha = (0,0)$
  \item $|\beta| = 1$: $\beta_1 = (1,0)$, $\beta_2 = (0,1)$
  \item $|\gamma| = 2$: $\gamma_1 = (1,1)$, $\gamma_2 = (2,0)$, $\gamma_3 = (0,2)$
  \item $|\delta| = 3$: $\delta_1 = (3,0)$, $\delta_2 = (2,1)$, $\delta_3 = (1,2)$, $\delta_4 = (0,3)$
\end{itemize}
Recall that
$$\partial^\alpha = \partial_1^{\alpha_1}\partial_2^{\alpha_2}\cdots\partial_n^{\alpha_n}, \qquad \text{e.g.\ } \partial^{\delta_2} = \partial^{(2,1)} = \partial_1^2\partial_2,$$
and
$$\alpha! = \alpha_1!\cdots\alpha_n!, \qquad \text{e.g.\ } \gamma_2! = (2,0)! = 2\cdot 1 = 2.$$
So:
$$
\begin{aligned}
f(\bar x + h) = &\ f(\bar x) && (\alpha) \\
&+ \partial_1 f(\bar x)h_1 + \partial_2 f(\bar x)h_2 && (\beta_1, \beta_2)\\
&+ \partial_1\partial_2 f(\bar x)h_1h_2 + \partial_1^2 f(\bar x)\frac{h_1^2}{2} + \partial_2^2 f(\bar x)\frac{h_2^2}{2} && (\gamma_1, \gamma_2, \gamma_3)\\
&+ \partial_1^3 f(\bar x)\frac{h_1^3}{6} + \partial_1^2\partial_2 f(\bar x)\frac{h_1^2h_2}{2} && (\delta_1, \delta_2)\\
&+ \partial_1\partial_2^2 f(\bar x)\frac{h_1h_2^2}{2} + \partial_2^3 f(\bar x)\frac{h_2^3}{6} && (\delta_3, \delta_4)\\
&+ o(|h|^3)
\end{aligned}
$$
\end{exercisesolution}

\begin{ainote}
Corsin writes the recall line as
$\partial^\alpha = \partial_{\alpha_1}\partial_{\alpha_2}\cdots\partial_{\alpha_n}$ with the
example $\partial^{(2,1)} = \partial_2\partial_1$. Read literally that is the wrong object --- a
multi-index $\alpha$ records \textit{how many times} each $\partial_i$ is applied, so
$\partial^{(2,1)} = \partial_1^2\partial_2$. His own expansion above uses
$\partial_1^2\partial_2 f(\bar x)\,\tfrac{h_1^2h_2}{2}$ for $\delta_2$, confirming the intended
meaning. Corrected above. See \texttt{OQ-10}.
\end{ainote}

We can rewrite this as
$$f(\bar x + h) = \sum_{\ell=0}^{k}\frac{1}{\ell!}\big(h_1\partial_1 + h_2\partial_2\big)^\ell f(\bar x).$$

\begin{theorem}[Taylor working formula]
\label{thm:taylor_working_formula}
For $f \in C^k(U \subseteq \mathbb{R}^n, \mathbb{R})$ with $\bar x + th \in U$ for all
$t \in [0,1]$:
$$f(\bar x + h) = \sum_{\ell=0}^{k}\frac{1}{\ell!}\left(\sum_{m=1}^{n} h_m\partial_m\right)^{\!\ell} f(\bar x) + o(|h|^k)$$
\end{theorem}

\begin{ainote}
Both sums in \cref{thm:taylor_working_formula} are written with lower limit $\ell = 1$ (resp.\
starting index) in the source. The $\ell = 0$ term is $f(\bar x)$ itself, which the explicit
expansion above does include, so the lower limit must be $0$. Corrected above. See
\texttt{OQ-11}.
\end{ainote}

\begin{exercise}[Taylor expansion of $(x+y^2)e^{-x^2-y^2}$]
\label{ex:taylor_explicit_example}
Compute the Taylor expansion of
$$f : \mathbb{R}^2\to\mathbb{R}, \qquad (x,y) \mapsto (x+y^2)e^{-x^2-y^2}$$
around $(0,0)$ up to third order.
\end{exercise}

\begin{exercisesolution}[\cref{ex:taylor_explicit_example}, cumbersome solution]
Calculate all partial derivatives up to third order and use the definition. (Please never do
this.)
\end{exercisesolution}

\begin{exercisesolution}[\cref{ex:taylor_explicit_example}, the easy solution]
The substitution and multiplication rules we used in Analysis 1 still apply:
$$
\begin{aligned}
(x+y^2)e^{-x^2-y^2} &= (x+y^2)\left(1 - x^2 - y^2 + o(|x^2+y^2|)\right) \\
&= x + y^2 - x^3 - xy^2 + o\!\left(|(x,y)|^3\right)
\end{aligned}
$$
\end{exercisesolution}

\begin{remark}[Second-order Taylor formula with Hessian]
The second-order Taylor approximation of $f \in C^2(\mathbb{R}^n,\mathbb{R})$ around
$x_0 \in \mathbb{R}^n$ can be expressed compactly as
$$f(x_0+x) = f(x_0) + \nabla f(x_0)\cdot x + \tfrac{1}{2}x\transp \cdot \Hess f(x_0)\cdot x + o(|x|^2)$$
where
$$\Hess f(x_0) = \big(\partial_i\partial_j f(x_0)\big)_{i,j=1,\dots,n}$$
is the \newterm{Hessian matrix} (\germanterm{Hesse-Matrix}) of $f$.
\end{remark}

\begin{question}
For $f \in C^k(U, \mathbb{R})$, do the Taylor expansions of $t \mapsto f(x_0+th)$ in $t$ and of
$(th) \mapsto f(x_0+th)$ in $(th)$ around $0$ coincide?
\end{question}

\begin{answer}
\textbf{Yes.} This is, in fact, how we prove the multi-dimensional Taylor approximation, and why
we need the assumption $x_0 + th \in U$ for all $t \in [0,1]$.
\end{answer}

\begin{exercise}[Proof of second-order Taylor formula]
\label{ex:hessian_taylor_proof}
Prove the formula from above:
$$f(x_0+h) = f(x_0) + \nabla f(x_0)h + \tfrac{1}{2}h\transp \Hess f(x_0)h + \dots$$
for $f \in C^k(U,\mathbb{R})$ with $x_0 + th \in U$ for all $t \in [0,1]$.
\end{exercise}

\begin{exercisesolution}[Proof of \cref{ex:hessian_taylor_proof}]
We expand $t \mapsto f(x_0+th)$ in $t$:
$$
\begin{aligned}
f(x_0+th) &= f(x_0) + t\left.\frac{d}{dt}\right|_{t=0}f(x_0+th) + \frac{t^2}{2}\left.\frac{d^2}{dt^2}\right|_{t=0}f(x_0+th) + o(|th|^2) \\
&= f(x_0) + t\sum_{i=1}^{n}\partial_i f(x_0)h_i + \frac{t^2}{2}\frac{d}{dt}\left.\sum_{i=1}^{n}\partial_i f(x_0+th)h_i\right|_{t=0} + o(|h|^2) \\
&= f(x_0) + t\nabla f(x_0)h + \frac{t^2}{2}\sum_{i,j=1}^{n}\partial_i\partial_j f(x_0)h_ih_j + o(|th|^2) \\
&= f(x_0) + t\nabla f(x_0)h + \frac{t^2}{2}h\transp \Hess f(x_0)h + o(|th|^2)
\end{aligned}
$$
Setting $t = 1$ gives the result.
\end{exercisesolution}
